\section{Introduction}

Detectors of machine-generated text report accuracy at or above 97\% on the benchmarks the field uses most, and frequently above 99\%~\cite{Guo2023close,Su2023plus,Zhou2024exploiting,Ansary2026multirepresentation}. Read directly, those numbers say the task is close to solved on the domains these corpora represent. Read carefully, they say something weaker, because a benchmark score measures separability of a particular corpus and not the capability a reader infers from it.

The difference matters because separability has more than one source. A detector can reach a high score by modelling how machine-generated language differs from human language, or by modelling how one corpus was assembled, punctuated and encoded. Both produce the same headline figure. Individual instances of the second have been documented, including a whitespace convention in HC3~\cite{Tian2023multiscale}, a length confound in the same corpus~\cite{Baidya2026detecting}, punctuation inconsistencies requiring standardisation~\cite{Ardeshirifar2025comparing}, and prompt-specific collection shortcuts~\cite{Park2024investigating}. What is missing is a way to ask how much, in total, a given benchmark's separability owes to surface form rather than content, and to compare benchmarks on that axis.

This paper proposes that measurement and applies it. We compare three arms on identical splits. A \emph{surface-only} model reads 47 orthographic features and never sees a word. A \emph{content-only} model sees text after punctuation, casing and non-ASCII characters have been stripped away. A \emph{full} model is a fine-tuned transformer on raw text. The gap between the first two bounds how much of the benchmark a detector could pass without reading language at all.

Applied to DAIGT V2 and HC3 across five model families, the answer differs sharply between the two. On HC3 the surface-only and content-only arms are indistinguishable, 3.20\% against 3.26\% test error. On DAIGT V2 content is stronger by a factor of 7.9 in error, 0.99\% against 7.86\%. The same contrast appears in the model comparison, where the best classical configuration on DAIGT V2 falls 0.0007 F1 short of the best transformer, inside our measured seed range, but falls short by an error-rate factor of 16 on HC3.

We also report two negative results, because both correct claims that this analysis originally made and that the literature invites. First, the whitespace cue that dominates discussion of HC3 is sufficient in isolation yet unnecessary in practice. BERT's tokeniser provably cannot represent it, emitting identical identifiers for strings that differ only in that respect, and BERT still reaches 0.9916 F1. Second, a character-level perturbation that drives DeBERTa from 0.9972 to 0.3644 F1 looks like a targeted vulnerability and is not one. The same model labels uniform random character strings as human in 100\% of cases at 0.98 confidence, so the collapse is not distinguishable from a degenerate response to unreadable input without a label-free control. We report that requirement rather than the vulnerability.

Our contributions are the decomposition itself, as a reusable measurement over any human-versus-machine corpus. A second contribution is its application to two widely used benchmarks, with a complete eight-configuration model comparison and per-model ROC, area and confusion analysis. A third is the set of three controls, on tokenisation, on pipeline execution and on label-free inputs, each of which overturns a conclusion that the uncontrolled version of this study had reached.
